\section{Inverse}
\bx{
\en{
    \item 
    Let $R_\theta$ be rotation counterclockwise by an angle $\theta$. 
    How would you express the inverse $R^{-1}_\theta$ as $R_\varphi$ for some $\varphi$ in a simple way? 
    If 
    \begin{equation*}
        A(\theta) = 
        \begin{pmatrix}
            \cos \theta & - \sin \theta \\
            \sin \theta & \cos \theta
        \end{pmatrix}
    \end{equation*}
    is the matrix associated with $R_\theta$, what is the matrix associated with $R^{-1}_\theta$? 
    \begin{Solution}
        $R^{-1}_\theta = R_{- \theta}$ because $R_{\theta} \circ R_{-\theta} = R_{\theta - \theta} = R_0 = I$. 
        The matrix associated with $R^{-1}_\theta$ is 
        \begin{equation*}
            \begin{pmatrix}
                \cos \theta & \sin \theta \\ 
                -\sin \theta & \cos \theta 
            \end{pmatrix}
        \end{equation*}
        because $\cos(-\theta) = \cos \theta$. 
    \end{Solution}
    \item 
    \ea{
        \item 
        Finish the proof of Theorem 2.1. 
        \item 
        Give the proof of Theorem 2.4 (ii). 
    }
    \begin{solution}
    \ea{
        \item  
        \begin{proof}
            Let $v \in V$. 
            Then we must show that 
            \begin{equation*}
                G(cv) = cG(v). 
            \end{equation*}
            Let $u = G(v)$. 
            By definition, this means that 
            \begin{equation*}
                F(u) = F(G(v)) = (F \circ G)(v) = I_V v = v. 
            \end{equation*}
            Since $F$ is linear, we find that 
            \begin{equation*}
                F(cu) = cF(u) = cv. 
            \end{equation*}
            By definition of the inverse map, this means that 
            \begin{equation*}
                G(cv) = G(F(cu)) = (G \circ F)(cu) = I_U(cu) = cv, 
            \end{equation*}
            thus proving what we wanted. 
        \end{proof}
        \item 
        \begin{proof}
            Suppose first that $F$ is surjective. 
            Then by the definition of surjection, we have 
            \begin{equation*}
                \Im F = W \quad \text{so} \quad \dim \Im F = \dim W. 
            \end{equation*}
            But
            \begin{equation*}
                \dim V = \dim \Ker F = \dim \Im F \quad \text{and} \quad \dim V = \dim W, 
            \end{equation*}
            so $\dim \Ker F = 0$, which implicates that $\Ker F = \{ O \}$. 
            Hence $F$ is injective by Theorem 2.2. 
            This proves (ii), by using Theorem 2.3. 
        \end{proof}
    }
    \end{solution}
    \item 
    \item 
    \item 
    \item 
    \item 
    \item 
    \item 
    \item 
    \item 
    Let $V$ be a two-dimensional vector space, and let $L: V \to V$ be a linear map such that $L^2 = O$, but $L \neq O$. 
    Let $v$ be an element of $V$ such that $L(v) \neq O$. 
    Let $w = L(v)$. 
    Prove that $\{ v, w \}$ is a basis of $V$. 
    \begin{Solution}
        \begin{proof}
            It suffices to prove that $v$ and $w$ are linearly independent. 
            Suppose $xv + yw = O$. 
            Apply $L$. 
            Then 
            \begin{equation*}
                \begin{aligned}
                    L(O) &= L(xv +yw) \\
                    &= xL(v) + yL(w) \\
                    &= xL(v) + yL(L(v)) \\
                    &= xL(v) + yO \\
                    &= O
                \end{aligned}
            \end{equation*}
            because $L^2 = O$. 
            Hence $xL(v) = O$. 
            Since $L(v) \neq O$, it follows that $x = 0$. 
            Then $y = 0$ because $w \neq O$. 
        \end{proof}
    \end{Solution}
    \item 
    \item 
}}